\chapter{PLANTEAMIENTO DEL PROBLEMA}
\thispagestyle{empty}

\section{DIAGNÓSTICO}
\label{sec:diagnostico}

En el Perú, entre 2010 y 2025, la máxima demanda del Sistema Eléctrico Interconectado Nacional (SEIN) aumentó de \SI{4579}{\mega\watt} a \SI{7942}{\mega\watt}. En el mismo período, la participación conjunta de la generación solar y eólica pasó de un valor cercano a cero a 9,5\,\pct{}. Estos cambios aumentan la exigencia de transportar energía y aprovechar la capacidad disponible de la red \parencites{minem2025anuario2024}{minem2025indicadoresDiciembre}{coesMaximaDemanda}{iea2025transmissiongrid}.

Los cables enterrados se usan en redes colectoras internas de las plantas solares y eólicas; así como, en tramos donde el espacio, la servidumbre, la seguridad o la integración urbana dificultan construir líneas aéreas \parencite{cigre2023overheadLines}. Su capacidad de transporte depende, además de las características propias del cable\footnote{La estructura del conductor y de las capas concéntricas del cable se describe en la Sección~\ref{subsec:generacion-evacuacion-calor} y se representa en la Figura~\ref{fig:cable-multicapa}.}, de la evacuación del calor generado por la corriente hacia el lecho de asiento, el suelo de relleno y el suelo nativo en el entorno del cable. La ampacidad \(\Imax\) es la corriente que el cable puede transportar de manera sostenida, bajo las condiciones de instalación especificadas, sin que la temperatura máxima del conductor \(\Tmax\) supere el límite térmico admisible \(T_{\mathrm{lim}}\) \parencites[cláusulas~4--5]{iec60287}[3--5]{enescu2020}.

Por tanto, el objeto de estudio es el arreglo de cables--entorno térmico, conformado por los cables, su disposición y los materiales circundantes que determinan conjuntamente la transferencia de calor \parencites[1--3]{enescu2020}[41--42]{cigre2025}. Este entorno es heterogéneo cuando su conductividad térmica \footnote{Debe notarse que algunas fuentes caracterizan el suelo mediante su resistividad térmica \(\rhoth\), que es la inversa de la conductividad, \(\rhoth=1/k\). La conductividad térmica $k$ se expresa en \si{\watt\per\metre\per\kelvin}; la resistividad térmica $\rhoth$ se expresa en \si{\kelvin\metre\per\watt}. Por tanto, una mayor resistividad equivale a una menor conductividad.} \(\kx\) varía espacialmente por la humedad, los estratos de suelos, los distintos suelos de relleno o las interfaces entre materiales. Esa variación cambia el campo de temperatura \(T(x,y)\), del cual se obtiene la temperatura máxima del conductor \(\Tmax\) y, mediante el límite térmico, la ampacidad \(\Imax\) \parencites[3--4]{enescu2021}[1,6--7]{khumalo2025}[3]{xing2023}.

El problema real aparece cuando una representación homogénea equivalente sustituye esa distribución espacial y se usa para estimar la ampacidad. La simplificación puede ocultar zonas de baja disipación y conducir a una estimación insegura o conservadora. En cambio, una descripción detallada mediante métodos de elementos finitos (FEM) o similares requieren construir, mallar y verificar cada escenario \parencites[41--42,93--99]{cigre2025}[1--3,12]{aldulaimi2024}.

El problema subyacente consiste en establecer cuándo y en qué magnitud la heterogeneidad cambia el campo de temperatura \(T(x,y)\), la temperatura máxima del conductor \(\Tmax\) y la ampacidad de los cables \(\Imax\) en comparación con la aproximación homogénea.

El problema técnico de la investigación consiste en diseñar y evaluar un artefacto basado en una red neuronal informada por física (PINN). La red incorporará la ecuación de conducción de calor y sus condiciones de frontera e interfaz para aproximar \(T(x,y)\) \parencite{raissi2019}. El artefacto integrará la representación de \(\kx\), la estimación de \(\Tmax\) e \(\Imax\), y los procedimientos de verificación.

La investigación sigue el paradigma de la ciencia del diseño (DSR) y adapta la secuencia de \textcite[54--56]{peffers2007}: identificación y motivación del problema, definición de los objetivos de solución, diseño y desarrollo, demostración, evaluación y comunicación. En este estudio, la evaluación corresponde a la verificación del artefacto; no se realizará validación experimental en campo. Esta adaptación conecta los tres niveles del problema con la construcción y verificación de la solución propuesta.

El primer capítulo desarrolla el diagnóstico, identifica el problema con sus antecedentes bibliográficos, lo formula y delimita su justificación y sus alcances. El segundo convierte esa formulación en objetivos, hipótesis y dimensiones operativas. El tercero reúne las bases teóricas y la definición de términos. El cuarto presenta la metodología DSR y la estrategia de verificación. El quinto administra el plan: cronograma, presupuesto y financiamiento. Las referencias y los anexos documentan las fuentes, la consistencia, la operacionalización, el diseño discreto de casos y la revisión documental.

\subsection{PROBLEMÁTICA}

Las redes eléctricas deben atender una demanda siempre creciente e integrar nuevas centrales de generación, especialmente renovables, ubicadas en distintos puntos del sistema. En América Latina y el Caribe, la generación eléctrica creció 5,5\,\pct{} en 2024 y las fuentes renovables aportaron 69\,\pct{}. De esta participación, 45\,\pct{} correspondieron a centrales hidroeléctricas convencionales y 19 a generación eólica y solar \parencite{olade2025generacionALC2024}.

En el Perú, entre 2010 y 2025, la producción eléctrica nacional pasó de \SI{35908}{\giga\watt\hour} a \SI{65109}{\giga\watt\hour}. El aumento acumulado fue 81\,\pct{} y la tasa media anual, 4\,\pct{}. En el mismo período, la máxima demanda del SEIN subió de \SI{4579}{\mega\watt} a \SI{7942}{\mega\watt}: 73\,\pct{} acumulado y 3,7\,\pct{} anualizado. Entre 2023 y 2025, la producción aumentó 4,9\,\pct{} y la máxima demanda, 4,4\,\pct{} \parencites[sec.~10.4]{minem2025anuario2024}[cuadro~4]{minem2025indicadoresDiciembre}{coesMaximaDemanda}.

La longitud acumulada de las líneas de 500 y \SI{220}{\kilo\volt} del SEIN aumentó de \SI{8266}{\kilo\metre} en 2010 a \SI{17179}{\kilo\metre} en 2025, equivalente a un 108\,\pct{} de crecimiento. La Figura~\ref{fig:demanda-expansion-red-peru} compara este inventario con la máxima demanda. Desde 2020, el inventario reportado presenta un crecimiento relativo menor que la demanda; lo que sugiere una mayor presión sobre las redes eléctricas \parencites{coesEstadisticasAnuales}{coesMaximaDemanda}.

\begin{figure}[htbp]
\centering
\begin{tikzpicture}
\begin{axis}[
  name=expansion,
  width=0.88\textwidth,
  height=0.52\textwidth,
  xmin=2009.5, xmax=2025.5,
  ymin=7500, ymax=20000,
  xtick={2010,2011,2012,2013,2014,2015,2016,2017,2018,2019,2020,2021,2022,2023,2024,2025},
  scaled x ticks=false,
  xticklabel={\pgfmathprintnumber[1000 sep={}]{\tick}},
  xticklabel style={rotate=45,anchor=east,font=\footnotesize},
  ytick={7500,10000,12500,15000,17500,20000},
  scaled y ticks=false,
  yticklabel={\pgfmathprintnumber[1000 sep={\,}]{\tick}},
  yticklabel style={font=\footnotesize,text=unired},
  xlabel={Año},
  ylabel={Líneas de 500 y 220 kV (km)},
  ylabel style={text=unired,font=\footnotesize},
  xlabel style={font=\footnotesize},
  axis y line*=right,
  axis x line*=bottom,
  grid=major,
  grid style={gray!20},
  legend style={at={(0.02,0.87)},anchor=north west,font=\scriptsize,draw=none,fill=white,fill opacity=0.88,text opacity=1}
]
\pgfplotsextra{
  \foreach \x/\lo/\hi in {
    2010/7500/8266, 2011/8266/9751, 2012/9751/10383,
    2013/10383/11569, 2014/11569/12573, 2015/12573/13453,
    2016/13453/14139, 2017/14139/15509, 2018/15509/15782,
    2019/15782/16510, 2020/16510/16383, 2021/16383/16101,
    2022/16101/16350, 2023/16350/16781, 2024/16781/17133,
    2025/17133/17179}{
      \pgfmathsetmacro{\xl}{\x-0.27}
      \pgfmathsetmacro{\xr}{\x+0.27}
      \draw[fill=unired!35,draw=unired!75]
        (axis cs:\xl,\lo) rectangle (axis cs:\xr,\hi);
  }
}
\addlegendimage{area legend,fill=unired!35,draw=unired!75}
\addlegendentry{Variación anual de la longitud acumulada}
\end{axis}

\begin{axis}[
  at={(expansion.south west)},anchor=south west,
  width=0.88\textwidth,
  height=0.52\textwidth,
  xmin=2009.5, xmax=2025.5,
  ymin=4400, ymax=8200,
  ytick={4500,5000,5500,6000,6500,7000,7500,8000},
  yticklabel style={font=\footnotesize,text=darkblue},
  ylabel={Máxima demanda del SEIN (MW)},
  ylabel style={text=darkblue,font=\footnotesize},
  axis y line*=left,
  axis x line=none,
  xtick=\empty,
  legend style={at={(0.02,0.97)},anchor=north west,font=\scriptsize,draw=none,fill=white,fill opacity=0.88,text opacity=1},
  clip=false
]
\addplot[color=darkblue,mark=*,mark size=2.0pt,very thick] coordinates {
  (2010,4578.94) (2011,4961.19) (2012,5291.00) (2013,5575.24)
  (2014,5737.27) (2015,6275.00) (2016,6492.41) (2017,6559.06)
  (2018,6884.59) (2019,7017.57) (2020,7125.30) (2021,7173.03)
  (2022,7467.45) (2023,7605.51) (2024,7794.01) (2025,7942.00)
};
\addlegendentry{Máxima demanda}
\end{axis}
\end{tikzpicture}
\caption[Máxima demanda y longitud de la red troncal]{Evolución de la máxima demanda del SEIN y de la longitud acumulada de líneas troncales de 500 y \SI{220}{\kilo\volt}, 2010--2025. Fuente: Elaboración propia con datos de las estadísticas anuales y del reporte de máxima demanda del COES \parencites{coesEstadisticasAnuales}{coesMaximaDemanda}.}
\label{fig:demanda-expansion-red-peru}
\end{figure}

La exigencia sobre la red también cambia con la composición y ubicación de la generación. En 2025, las fuentes renovables aportaron aproximadamente 63\,\pct{} de la producción nacional, donde 52\,\pct{} correspondieron a centrales hidroeléctricas convencionales y 11 a recursos no convencionales. Como muestra la Figura~\ref{fig:participacion-renovable-peru}, la participación eólica y solar aumentó de un valor cercano a cero en 2010 a 9,5\,\pct{} en 2025. Entre 2023 y 2025 creció 4,2 puntos \parencites[sec.~10.4]{minem2025anuario2024}[cuadro~4]{minem2025indicadoresDiciembre}. Esta generación requiere redes colectoras, conexiones y capacidad de transporte entre nuevos puntos del SEIN \parencites{coes2023diagnostico2025_2034}{iea2025transmissiongrid}.

\begin{figure}[htbp]
\centering
\begin{tikzpicture}
\begin{axis}[
  name=variacion,
  width=0.88\textwidth,
  height=0.55\textwidth,
  ylabel={Variación interanual (\pct{})},
  xmin=2009.5, xmax=2025.5,
  ymin=-0.5, ymax=3.1,
  ytick={-0.5,0,0.5,1,1.5,2,2.5,3},
  axis y line*=right,
  axis x line=none,
  xtick=\empty,
  ybar,
  bar width=7pt,
  tick label style={font=\footnotesize},
  label style={font=\footnotesize}
]
\addplot[fill=lightblue,draw=darkblue!60] coordinates {
  (2011,0.000)
  (2012,0.135)
  (2013,0.319)
  (2014,0.546)
  (2015,0.708)
  (2016,0.813)
  (2017,0.058)
  (2018,1.513)
  (2019,0.150)
  (2020,0.670)
  (2021,-0.343)
  (2022,0.038)
  (2023,0.726)
  (2024,2.788)
  (2025,1.386)
};
\end{axis}
\begin{axis}[
  at={(variacion.south west)},
  anchor=south west,
  width=0.88\textwidth,
  height=0.55\textwidth,
  ylabel={Participación en la generación nacional (\pct{})},
  xlabel={Año},
  xmin=2009.5, xmax=2025.5,
  xtick={2010,2012,2014,2016,2018,2020,2022,2024,2025},
  xticklabel={\pgfmathprintnumber[fixed,precision=0,1000 sep={}]{\tick}},
  xticklabel style={rotate=45,anchor=east,font=\footnotesize},
  ymin=0, ymax=11,
  ytick={0,2,4,6,8,10},
  axis y line*=left,
  axis x line*=bottom,
  grid=major,
  grid style={gray!25},
  axis line style={black!70,line width=0.8pt},
  tick label style={font=\footnotesize},
  label style={font=\footnotesize},
  legend style={font=\footnotesize,at={(0.02,0.98)},anchor=north west,draw=none,fill=white}
]
\addplot[color=darkblue,mark=*,mark size=2.2pt,very thick] coordinates {
  (2010,0.003)
  (2011,0.003)
  (2012,0.138)
  (2013,0.457)
  (2014,1.003)
  (2015,1.711)
  (2016,2.524)
  (2017,2.582)
  (2018,4.095)
  (2019,4.245)
  (2020,4.915)
  (2021,4.572)
  (2022,4.610)
  (2023,5.336)
  (2024,8.124)
  (2025,9.510)
};
\addlegendentry{Participación solar + eólica}
\addlegendimage{area legend,fill=lightblue,draw=darkblue!60}
\addlegendentry{Variación interanual}
\end{axis}
\end{tikzpicture}
\caption[Participación de la generación solar y eólica en el Perú]{Participación conjunta de la generación solar y eólica en la producción eléctrica nacional del Perú y variación interanual, 2010--2025. Fuente: Elaboración propia con datos del Ministerio de Energía y Minas \parencites[sec.~10.4]{minem2025anuario2024}[cuadro~4]{minem2025indicadoresDiciembre}. Notas: las barras expresan diferencias en puntos porcentuales, no tasas relativas; el valor de 2025 corresponde a cifras preliminares.}
\label{fig:participacion-renovable-peru}
\end{figure}

Los cables enterrados forman parte de las redes colectoras internas de las plantas eólicas y solares. Entre 2023 y 2025 se conectaron al SEIN siete centrales de al menos \SI{80}{\mega\watt} ---cuatro eólicas y tres solares---, con aproximadamente \SI{1250}{\mega\watt} en conjunto \parencites{osinergmin2024puntaLomitas}{minem2024nuevasRER}{minem2025nuevasCentrales}. La Central Eólica Punta Lomitas muestra la escala de esta infraestructura, pues sus 57 aerogeneradores se conectan mediante circuitos colectores de \SI{33}{\kilo\volt} constituidos por \SI{290}{\kilo\metre} de cables unipolares XLPE\footnote{XLPE corresponde a \textit{cross-linked polyethylene} o polietileno reticulado, un material aislante termoestable empleado en cables de potencia \parencites[9--10]{thue1999electrical}{iec60502_2_2014}. La estructura multicapa de estos cables se describe en la Sección~\ref{subsec:generacion-evacuacion-calor} y se muestra en la Figura~\ref{fig:cable-multicapa}.} \parencite{coes2021puntaLomitasEO}.

El uso de cables enterrados no se limita a la generación renovable. También aparece cuando el espacio, la servidumbre, la seguridad, los cruces o la integración urbana dificultan construir líneas aéreas \parencite{cigre2023overheadLines}. En Lima, la reconfiguración Chavarría--Santa Rosa--Carapongo incluye un circuito subterráneo de \SI{220}{\kilo\volt} \parencite[13]{coes2022planTransmision2023_2032}. En Huancayo, el proyecto Palián contempla aproximadamente \SI{10}{\kilo\metre} de trazados subterráneos con cables de \SI{60}{\kilo\volt} \parencites[168]{coes2022planTransmision2023_2032}[3,55--56]{proinversion2026palian}.

En conjunto, estos casos muestran que el sistema eléctrico peruano incorpora cables enterrados tanto en redes colectoras de generación como en enlaces de transmisión. La extensión y función de esta infraestructura hacen relevante estimar adecuadamente su ampacidad para sustentar decisiones de diseño y operación térmica, evitar restricciones innecesarias de capacidad y contribuir a la operación confiable de los activos conectados.

El dimensionamiento y la operación de los cables requieren estimar su ampacidad: la corriente máxima que pueden transportar sin superar su temperatura admisible \parencite{iec60287}. Este límite depende del cable y de la disipación de calor a través del lecho de asiento, el suelo de relleno y el suelo nativo \parencites{fariz2026}{mobius2025}. Por tanto, el análisis debe considerar el arreglo y su entorno enterrado como un sistema.

Los métodos clásicos y las normas aportan la base de cálculo, pero simplifican el entorno mediante resistencias térmicas equivalentes y propiedades homogéneas \parencites[752--760]{neher1957}{iec60287}{iec60853}. Esta representación puede perder detalle cuando la humedad, los estratos, el lecho de asiento, el suelo de relleno o las zonas secas cambian la conductividad térmica alrededor del cable \parencites[4--6]{enescu2021}[1--3]{aldulaimi2024}. Por ello, instalaciones con cables y corriente iguales pueden alcanzar temperaturas distintas si cambia el entorno térmico.

El impacto de las características térmicas del entorno sobre la ampacidad está documentado por varios autores. En Sandton, Sudáfrica, \textcite[5--6, tabs.~2--3]{khumalo2025} analizaron un cable tripolar XLPE de media tensión en configuración $1\times3\times\SI{185}{\milli\metre\squared}$, con conductores de cobre y entierro directo en suelo nativo. El aumento de la resistividad térmica de \SI{0.6}{\kelvin\metre\per\watt} a \SI{3.7}{\kelvin\metre\per\watt} redujo la ampacidad de \SI{518}{\ampere} a \SI{224}{\ampere}.

En una planta fotovoltaica ubicada en el sur del Perú, \textcite[7--8]{ariasVelasquez2023} estudiaron tres cables unipolares XLPE de aluminio de \SI{630}{\milli\metre\squared}, en disposición trébol y a \SI{33}{\kilo\volt}. Obtuvieron \SI{563}{\ampere} para un suelo uniforme de \SI{1.5}{\kelvin\metre\per\watt}, frente a \SI{350}{\ampere} con la resistividad medida de \SI{4.9}{\kelvin\metre\per\watt}. Esta comparación también incorpora las condiciones reales de temperatura del terreno.

En un caso simulado sin ubicación geográfica reportada, \textcite[15--16]{atoccsa2024} analizaron tres cables unipolares XLPE de \SI{220}{\kilo\volt} y \SI{2000}{\milli\metre\squared}, dispuestos horizontalmente dentro de un suelo de relleno con mejora térmica. Estimaron que elevar la resistividad del suelo de \SI{2.5}{\kelvin\metre\per\watt} a \SI{3.0}{\kelvin\metre\per\watt} redujo la ampacidad de \SI{1157}{\ampere} a \SI{1064}{\ampere}.

Por último, \textcite[12--16]{kim2025} modelaron en Corea del Sur un banco de seis cables XLPE de \SI{154}{\kilo\volt}, con conductores de cobre de \SI{1200}{\milli\metre\squared}, y compararon el uso de hormigón de agregado precolocado (PAC, $k=\SI{2.1}{\watt\per\metre\per\kelvin}$) con arena natural ($k=\SI{1.4}{\watt\per\metre\per\kelvin}$). El PAC elevó la ampacidad de aproximadamente \SI{941}{\ampere} a \SI{988}{\ampere} y, a igual corriente, redujo la temperatura máxima de \SI{78}{\celsius} a \SI{71}{\celsius}. La Figura~\ref{fig:resistividad-ampacidad-casos} reúne los cuatro casos.

\begin{figure}[htbp]
\centering
\begin{tikzpicture}
\begin{groupplot}[
  group style={group size=4 by 1,horizontal sep=0.35cm},
  width=0.30\textwidth,
  height=0.54\textwidth,
  ymin=0, ymax=1600,
  xmin=0.05, xmax=0.95,
  xtick=\empty,
  ytick={0,200,400,600,800,1000,1200},
  grid=major,
  grid style={gray!25},
  axis line style={black!70},
  tick label style={font=\scriptsize},
  label style={font=\scriptsize},
  xlabel style={font=\tiny,align=center},
  clip=false
]
% Cada panel conserva un color; el tono claro representa la menor
% resistividad y el oscuro la mayor resistividad.
\nextgroupplot[
  ylabel={Ampacidad (\si{\ampere})},
  xlabel={\shortstack{Khumalo et al.\\MT; $1\times3\times\SI{185}{\milli\metre\squared}$\\Cu; tripolar\\entierro directo}}
]
\addplot[ybar,bar width=16pt,fill=darkblue!35,draw=darkblue] coordinates {(0.32,518)};
\addplot[ybar,bar width=16pt,fill=darkblue!85,draw=darkblue] coordinates {(0.68,224)};
\node[font=\scriptsize,align=center,anchor=north] at (axis cs:0.50,1580) {\textcolor{darkblue!55}{\rule{1.1ex}{1.1ex}} $\rhoth=0{,}6$\\$k=1{,}7$; 518 A};
\node[font=\scriptsize,align=center,anchor=north] at (axis cs:0.50,1350) {\textcolor{darkblue!90}{\rule{1.1ex}{1.1ex}} $\rhoth=3{,}7$\\$k=0{,}3$; 224 A\\$-57$\,\pct{} $I_{\max}$};

\nextgroupplot[
  yticklabels=\empty,
  xlabel={\shortstack{Arias Velásquez et al.\\\SI{33}{\kilo\volt}; $3\times\SI{630}{\milli\metre\squared}$\\Al; unipolares\\trébol}}
]
\addplot[ybar,bar width=16pt,fill=unired!30,draw=unired] coordinates {(0.32,563)};
\addplot[ybar,bar width=16pt,fill=unired!80,draw=unired] coordinates {(0.68,350)};
\node[font=\scriptsize,align=center,anchor=north] at (axis cs:0.50,1580) {\textcolor{unired!45}{\rule{1.1ex}{1.1ex}} $\rhoth=1{,}5$\\$k=0{,}7$; 563 A};
\node[font=\scriptsize,align=center,anchor=north] at (axis cs:0.50,1350) {\textcolor{unired!90}{\rule{1.1ex}{1.1ex}} $\rhoth=4{,}9$\\$k=0{,}2$; 350 A\\$-38$\,\pct{} $I_{\max}$};

\nextgroupplot[
  yticklabels=\empty,
  xlabel={\shortstack{Atoccsa et al.\\\SI{220}{\kilo\volt}; $3\times\SI{2000}{\milli\metre\squared}$\\Cu; unipolares\\plana}}
]
\addplot[ybar,bar width=16pt,fill=teal!30,draw=teal!80!black] coordinates {(0.32,1157)};
\addplot[ybar,bar width=16pt,fill=teal!80,draw=teal!80!black] coordinates {(0.68,1064)};
\node[font=\scriptsize,align=center,anchor=north] at (axis cs:0.50,1580) {\textcolor{teal!45}{\rule{1.1ex}{1.1ex}} $\rhoth=2{,}5$\\$k=0{,}4$; 1157 A};
\node[font=\scriptsize,align=center,anchor=north] at (axis cs:0.50,1420) {\textcolor{teal!90!black}{\rule{1.1ex}{1.1ex}} $\rhoth=3{,}0$\\$k=0{,}3$; 1064 A$^{*}$\\$-8$\,\pct{} $I_{\max}$};

\nextgroupplot[
  yticklabels=\empty,
  xlabel={\shortstack{Kim et al.\\\SI{154}{\kilo\volt}; $6\times\SI{1200}{\milli\metre\squared}$\\Cu; unipolares\\2 filas planas}}
]
\addplot[ybar,bar width=16pt,fill=violet!30,draw=violet!80!black] coordinates {(0.32,988)};
\addplot[ybar,bar width=16pt,fill=violet!80,draw=violet!80!black] coordinates {(0.68,941)};
\node[font=\scriptsize,align=center,anchor=north] at (axis cs:0.50,1580) {\textcolor{violet!45}{\rule{1.1ex}{1.1ex}} $\rhoth=0{,}48$\\$k=2{,}1$; 988 A};
\node[font=\scriptsize,align=center,anchor=north] at (axis cs:0.50,1350) {\textcolor{violet!90!black}{\rule{1.1ex}{1.1ex}} $\rhoth=0{,}73$\\$k=1{,}4$; 941 A$^{\dagger}$\\$-5$\,\pct{} $I_{\max}$};
\end{groupplot}
\end{tikzpicture}
\caption[Resistividad térmica y ampacidad en cuatro casos publicados]{Cambio de ampacidad entre dos valores de resistividad térmica en cuatro instalaciones diferentes. Nota: la conductividad térmica es la inversa de la resistividad, $k=1/\rhoth$; en las etiquetas, $\rhoth$ se expresa en \si{\kelvin\metre\per\watt} y $k$ en \si{\watt\per\metre\per\kelvin}. $^{*}$Aproximación local de sensibilidad. $^{\dagger}$Valor leído aproximadamente de la Figura~15 de la fuente, para septiembre y una temperatura del terreno de \SI{21}{\celsius}. Fuente: Elaboración propia con datos de \textcite[5--6, tabs.~2--3]{khumalo2025}, \textcite[7--8]{ariasVelasquez2023}, \textcite[15--16]{atoccsa2024} y \textcite[4,12--13, fig.~15]{kim2025}.}
\label{fig:resistividad-ampacidad-casos}
\end{figure}


En conjunto, los casos muestran que las propiedades térmicas del suelo nativo y de los suelos de relleno modifican la ampacidad y el margen térmico disponible, y que además pueden cambiar durante la vida útil. Por ello, la estimación de ampacidad requiere supuestos documentados, datos del sitio y modelos que puedan verificarse y actualizarse \parencites{fariz2026}{mobius2025}.

La dificultad metodológica consiste en elegir una representación adecuada del entorno. El modelo homogéneo sustituye la distribución espacial por una propiedad única o promedio. Es útil como referencia, pero puede ocultar zonas locales de baja disipación. FEM representa la geometría y las propiedades espaciales, aunque requiere construir, mallar y verificar cada escenario \parencites[3--4,93--99]{cigre2025}[1--3]{aldulaimi2024}. La literatura revisada no ofrece, para el dominio delimitado, una evaluación reproducible que separe los efectos del contraste, el patrón, la extensión y la proximidad sobre $T(x,y)$, $\Tmax$ e $\Imax$.

La Figura~\ref{fig:brecha} ordena la cadena que conecta la presión de red con el riesgo de decisión. La incertidumbre no proviene solo del terreno, sino también de cómo se representa en el cálculo. Por ello, el estudio separará el efecto físico de la heterogeneidad del error numérico del artefacto. El primero se evaluará mediante escenarios controlados; el segundo, mediante referencias analíticas, manufacturadas, FEM y casos publicados \parencites[4--6]{enescu2021}[3--4,93--99]{cigre2025}.

Si el modelo subestima la temperatura, puede admitir una corriente excesiva y acelerar el envejecimiento del cable \parencite[1--2,9]{khumalo2025}. Si la sobreestima, puede restringir capacidad disponible y favorecer el sobredimensionamiento \parencites[4--6]{enescu2021}[1--2,9]{khumalo2025}. La consecuencia depende tanto de la representación física como de la exactitud del cálculo.

\vspace{0.15cm}

\begin{figure}[htbp]
\centering
\resizebox{1.0\textwidth}{!}{%
\begin{tikzpicture}[
  scale=1.,
  transform shape,
  box/.style={draw,rounded corners=3pt,align=center,text width=3.25cm,minimum height=1.25cm,font=\tiny},
  arr/.style={-{Stealth[length=3mm]},thick},
  darr/.style={-{Stealth[length=3mm]},thick,dashed}]
\node[box,fill=lightgray] (cables) {Tramos con cables\\enterrados\\restricciones de espacio,\\servidumbre y seguridad};
\node[box,fill=lightblue,right=0.75cm of cables] (real) {Entorno térmico\\enterrado\\$\kx$ variable: humedad,\\estratos, secado, suelo de relleno};
\node[box,fill=red!10,right=0.75cm of real] (gap) {Efecto no\\cuantificado sobre\\$T(x,y)$, $\Delta \Tmax$ y $\Delta \Imax$};
\node[box,fill=yellow!15,below=1.0cm of cables] (context) {Presión de red\\SEIN, demanda, renovables\\continuidad y capacidad};
\node[box,fill=green!10,below=1.0cm of real] (simpl) {Representación\\homogénea\\$k=k_{eq}$ y resistencias\\térmicas equivalentes};
\node[box,fill=orange!12,below=1.0cm of gap] (impact) {Riesgo de decisión\\subestimar: envejecimiento\\sobreestimar: capacidad ociosa};
\draw[arr] (context) -- (cables);
\draw[arr] (cables) -- (real);
\draw[arr] (real) -- (gap);
\draw[arr] (real) -- (simpl);
\draw[arr] (simpl) -- (gap);
\draw[darr] (context.south) -- ($(context.south)+(0,-0.55cm)$)
  -- ($(impact.south)+(0,-0.55cm)$) -- (impact.south);
\draw[arr] (gap) -- (impact);
\draw[arr] (simpl) -- (impact);
\end{tikzpicture}
}%
\caption{Relación entre presión de red, uso de cables enterrados, entorno térmico heterogéneo, simplificación homogénea y riesgo de decisión. Fuente: Elaboración propia.}
\label{fig:brecha}
\end{figure}

En resumen, la expansión y transformación del SEIN aumentan la exposición a instalaciones enterradas cuya capacidad depende de condiciones térmicas difíciles de observar. Una representación insuficiente puede producir decisiones inseguras o conservadoras. La brecha consiste en determinar, dentro del dominio 2D estacionario, cuándo y cuánto cambian $T(x,y)$, $\Tmax$ e $\Imax$ por el contraste, el patrón, la extensión y la proximidad de la heterogeneidad. Esta brecha conduce al planteamiento formal del problema y a la evaluación del artefacto PINN.

\section{IDENTIFICACIÓN Y DESCRIPCIÓN DEL PROBLEMA DE ESTUDIO}
\label{sec:planteamiento-problema}

El problema real se manifiesta cuando una estimación térmica que no representa adecuadamente el entorno enterrado sustenta decisiones de ampacidad. Si la temperatura se subestima, el cable puede operar por encima de su límite térmico; si se sobreestima, puede restringirse capacidad utilizable. La tesis no resolverá por sí sola la planificación ni la operación de la red, pero producirá evidencia para reducir esta incertidumbre dentro del dominio estudiado.

El problema subyacente es un vacío de conocimiento: no se ha establecido de forma reproducible cómo el contraste, el patrón, la extensión y la proximidad de las heterogeneidades de $\kx$ modifican el campo $T(x,y)$, $\Tmax$ y $\Imax$ frente a una representación homogénea equivalente, ni bajo qué condiciones esas diferencias son relevantes. Un promedio simple no responde necesariamente esta cuestión porque la geometría, las interfaces y las discontinuidades del suelo interactúan en la conducción del calor \parencites[4--6]{enescu2021}[3--4,93--99]{cigre2025}.

El problema técnico de investigación consiste en diseñar y evaluar una formulación computacional integrada, reproducible y verificada que represente explícitamente $\kx$, aproxime $T(x,y)$ y permita cuantificar su efecto sobre $\Tmax$ e $\Imax$. El artefacto PINN deberá documentar exactitud, consistencia física y robustez mediante referencias analíticas, manufacturadas, FEM y casos publicados, sin asumir que supera a los métodos normativos o numéricos existentes \parencites{iec60287}{iec60853}[3--4,93--99]{cigre2025}[686--690]{raissi2019}.

El objeto de estudio es el arreglo de cables--entorno térmico en instalaciones subterráneas. Es un objeto físico tangible formado por los cables multicapa y su disposición, el lecho de asiento, el suelo de relleno y el suelo nativo; su respuesta está determinada por las condiciones térmicas de operación \parencites[3--6]{enescu2021}{ieee442}[1,10]{kim2025}. La unidad de análisis será una sección transversal 2D de este objeto con cables de tipo XLPE, suelo nativo, suelo de relleno y, cuando corresponda, suelo de relleno con propiedades térmicas mejoradas.

El interés no recae en una marca específica de cable. En cambio, se centra en configuraciones representativas y reproducibles del objeto, tanto en disposición plana como en trébol. El artefacto PINN 2D será el producto tecnológico diseñado y evaluado para representarlo computacionalmente; por tanto, no se confundirá con el objeto físico de estudio.

\subsection{ANTECEDENTES BIBLIOGRÁFICOS}
\label{sec:antecedentes}

La Tabla~\ref{tab:antecedentes-cronologia} ordena los antecedentes de la fuente más antigua a la más reciente; los apartados siguientes los agrupan por tema y también siguen ese orden dentro de cada tema, salvo las revisiones, que se citan junto con el trabajo al que se refieren. Ocho de las 19 fuentes son anteriores a 2021 y se conservan por ser clásicos (\textcite{neher1957}), normas de la Comisión Electrotécnica Internacional (\textit{International Electrotechnical Commission}, IEC) \parencite{iec60853}, fundamentos metodológicos (\textcite{peffers2007}; \textcite{gregor2013}), el enfoque PINN original (\textcite{raissi2019}) o estudios y revisiones que siguen siendo referencia del tema (\textcite{aras2005}; \textcite{oclon2015}; \textcite{enescu2020}). Las 11 restantes son de 2021 o posteriores.

\begingroup\footnotesize\setlength{\tabcolsep}{3pt}
\begin{longtable}{L{1.1cm}L{3.4cm}L{7.6cm}L{2.4cm}}
\caption{Evolución cronológica de los antecedentes, de la fuente más antigua a la más reciente. Fuente: Elaboración propia a partir de las referencias citadas.}\label{tab:antecedentes-cronologia}\\
\toprule
Año & Fuente & Aporte al problema & Naturaleza \\
\midrule\endfirsthead
\toprule
Año & Fuente & Aporte al problema & Naturaleza \\
\midrule\endhead
1957 & \textcite{neher1957} & Formulación clásica de pérdidas, resistencias térmicas y elevación de temperatura. & Clásico \\
2002 & \textcite{iec60853} & Corriente cíclica y de emergencia: ciclos de carga y secado parcial. & Norma \\
2005 & \textcite{aras2005} & Comparación de enfoques analíticos y numéricos para cables enterrados; FEM representa fuentes cercanas. & Estudio de referencia \\
2007 & \textcite{peffers2007} & Seis actividades de la ciencia del diseño. & Método \\
2013 & \textcite{gregor2013} & Artefacto situado y contribuciones abstractas; evaluación por validez, utilidad, calidad y eficacia. & Método \\
2015 & \textcite{oclon2015} & Optimización del lecho de asiento: una región de alta conductividad mejora la disipación. & Estudio técnico \\
2019 & \textcite{raissi2019} & Enfoque PINN para problemas directos e inversos. & Fundamento PINN \\
2020 & \textcite{enescu2020} & Revisión de la evaluación térmica y del rating; métodos analíticos, circuitales y de campo. & Revisión \\
2021 & \textcite{enescu2021} & Revisión de conceptos y métodos de rating térmico dinámico; influencia del suelo y del secado. & Revisión \\
2022 & \textcite{lawal2022} & Revisión sistemática de PINN: ventajas y dificultades de entrenamiento. & Revisión \\
2023 & \textcite{xing2023} & PINN para conducción de calor anisótropa estacionaria 3D. & Estudio técnico \\
2023 & \textcite{iec60287} & Ecuaciones de corriente admisible y pérdidas en régimen permanente. & Norma \\
2024 & \textcite{atoccsa2024} & Optimización de ampacidad con relleno térmico en cables de alta tensión. & Estudio técnico \\
2024 & \textcite{delport2024} & Ciclos de evaluación durante el desarrollo; criterios de abstracción, originalidad, justificación y beneficio. & Método \\
2025 & \textcite{khumalo2025} & Relación entre humedad, resistividad, profundidad y temperatura del suelo. & Estudio técnico \\
2025 & \textcite{kim2025} & Reducción de temperatura con un material de relleno de mayor conductividad. & Estudio técnico \\
2025 & \textcite{pan2025} & Reconstrucción de un campo térmico 2D y de sus condiciones de borde con datos limitados. & Estudio técnico \\
2025 & \textcite{ren2025} & Revisión de la evolución metodológica de las PINN y de sus dificultades de entrenamiento. & Revisión \\
2025 & \textcite{cigre2025} & Casos heterogéneos y verificación de las herramientas de rating por elementos finitos. & Recomendación técnica \\
\bottomrule
\end{longtable}\endgroup

\subsubsection{MODELOS TÉRMICOS Y AMPACIDAD}

\textcite[752--772]{neher1957} establecieron una formulación clásica para relacionar pérdidas, resistencias térmicas y elevación de temperatura. Por tanto, su trabajo es una base histórica de los métodos de \textit{rating}.

La IEC 60853, de 2002, amplía el tratamiento a ciclos de carga y secado parcial \parencite[cláusulas 2--4]{iec60853}. \textcite[1387--1400]{aras2005} compararon enfoques analíticos y numéricos para cables enterrados; sus resultados muestran que FEM representa fuentes cercanas y distribuciones de temperatura mediante geometría explícita.

Las revisiones de Enescu y colaboradores ordenan los métodos analíticos, circuitales y de campo \parencite[4,11,17]{enescu2020}. Además, señalan que la conductividad del suelo, las capas y el secado son factores dominantes. La revisión sobre capacidad térmica dinámica añade una relación directa para el estudio: la resistividad térmica del suelo influye en la ampacidad, y un valor mayor la reduce \parencite[3--6]{enescu2021}.

La IEC 60287, en su edición de 2023, formaliza ecuaciones de corriente admisible y pérdidas para régimen permanente \parencite[cláusulas 4--5]{iec60287}. Además, las recomendaciones del Consejo Internacional de Grandes Redes Eléctricas (CIGRÉ) de 2025 documentan casos heterogéneos y exigen verificación antes de confiar en una herramienta de \textit{rating} \parencite[3--4,41--42,93--99]{cigre2025}. Por ello, la verificación será un criterio central del estudio.

Estos métodos son referencias necesarias para comparación. Sin embargo, no sustituyen una solución de campo cuando la topología térmica no cumple sus supuestos \parencites{iec60287}{iec60853}[3--4,93--99]{cigre2025}.

\subsubsection{HETEROGENEIDAD, SUELO NATIVO Y SUELO DE RELLENO}

Los trabajos de optimización del lecho de asiento y del suelo de relleno muestran que una región de alta conductividad puede mejorar la disipación \parencites[1--3]{oclon2015}[1--4]{atoccsa2024}. Sin embargo, se apoyan principalmente en FEM y optimización externa.

\textcite[1,6--9]{khumalo2025} aportan evidencia sobre humedad, resistividad, profundidad y temperatura ambiente del suelo. Por su parte, \textcite[1,10,16]{kim2025} cuantifican la reducción de temperatura lograda con un suelo de relleno más conductor. En conjunto, ambos trabajos muestran que $k$ no es una constante secundaria. Por tanto, debe tratarse como variable de diseño e incertidumbre. El aporte previsto de este estudio es distinto, porque busca trasladar la física a una PINN y evaluar reglas sobre contraste, posición y extensión de regiones heterogéneas.

\subsubsection{INTELIGENCIA ARTIFICIAL INFORMADA POR LA FÍSICA}

\textcite[686--690]{raissi2019} consolidaron el enfoque PINN para problemas directos e inversos al combinar datos y residuos de ecuaciones diferenciales. Las revisiones posteriores confirman su expansión, pero también identifican entrenamiento lento, competencia entre términos de pérdida, dificultad multiescala y necesidad de estrategias adaptativas o descomposición de dominios \parencites[1--2,14--17]{lawal2022}[2--6]{ren2025}.

En conducción de calor, \textcite[1,3,17]{xing2023} demostraron una PINN para un problema anisotrópico estacionario 3D y reportaron, en uno de sus ejemplos, error relativo de 0,23\,\pct. \textcite[1,3,16]{pan2025} reconstruyeron un campo térmico 2D y condiciones de borde con datos limitados, con error absoluto inferior a \SI{1}{\celsius} en el dominio de validación declarado. Estos resultados prueban viabilidad en problemas térmicos controlados, pero no eliminan la necesidad de evaluar interfaces, pérdidas internas y geometrías de cables.

\subsubsection{ANTECEDENTES METODOLÓGICOS DSR}

\textcite[54--56]{peffers2007} proponen seis actividades: identificar y motivar el problema, definir objetivos de solución, diseñar y desarrollar, demostrar, evaluar y comunicar. La secuencia no es rígida y permite iterar cuando la evidencia de evaluación no alcanza los objetivos.

\textcite[342--350]{gregor2013} distinguen el artefacto situado de las contribuciones más abstractas, tales como métodos y principios de diseño. También proponen evaluar validez, utilidad, calidad y eficacia, y comunicar con claridad el problema, el artefacto y la evidencia. \textcite[198--203]{delport2024} complementan esta visión con ciclos continuos de evaluación durante el desarrollo y con los principios de abstracción, originalidad, justificación y beneficio.

\subsubsection{UBICACIÓN DEL PROYECTO}

El proyecto se ubica en la intersección de tres líneas de trabajo: el cálculo térmico y la ampacidad de cables enterrados, las redes neuronales informadas por física y la ciencia del diseño de artefactos. Se desarrollará como tesis profesionalizante del programa de Maestría en Inteligencia Artificial de la Unidad de Posgrado de la Facultad de Ingeniería Industrial y de Sistemas de la Universidad Nacional de Ingeniería.

Respecto de la primera línea, no reemplaza las normas ni la referencia FEM: usará la IEC 60287 como línea base cuando el caso sea comparable y FEM como referencia independiente del campo. Respecto de la segunda, aportará una comparación trazable de una formulación PINN por material, con transmisión de temperatura y de flujo en las interfaces, aplicada a configuraciones de cable enterrado. Respecto de la tercera, corresponde a un artefacto situado con principios de diseño derivados, en el sentido de \textcite[342--350]{gregor2013}. La brecha descrita en la Sección~\ref{sec:planteamiento-problema} delimita el aporte: la comparación específica entre esas tres líneas, y no una novedad general sobre ninguna de ellas.

\subsection{FORMULACIÓN DEL PROBLEMA}

La formulación se presenta a continuación y se organiza en la matriz de consistencia del \ref{anx:matriz-consistencia}, la cual funciona como un organizador lógico del plan. Por tanto, la matriz relaciona problema, objetivos, hipótesis, variables, evidencia y controles, además de evitar que las partes del estudio queden como ideas separadas. Los problemas específicos siguen la lógica DSR. Por ello, delimitan entradas del problema, productos verificables y criterios de evaluación del artefacto.

\subsubsection{PROBLEMA GENERAL}

En el arreglo de cables--entorno térmico, la heterogeneidad espacial del suelo nativo, el lecho de asiento y el suelo de relleno puede modificar el campo de temperatura $T(x,y)$, la temperatura máxima $\Tmax$ y la ampacidad $\Imax$ respecto de una representación homogénea equivalente \parencites[4--6]{enescu2021}[3--4,93--99]{cigre2025}. Sin embargo, resulta difícil cuantificar esas diferencias de forma reproducible y establecer cuándo son relevantes para las decisiones de diseño y operación. Las representaciones homogéneas equivalentes son rápidas, aunque no describen de forma explícita la conductividad variable $\kx$; los modelos numéricos de elementos finitos (FEM) y similares sí pueden representarla, pero exigen construir, mallar y verificar cada escenario. En consecuencia, persiste la necesidad de una formulación computacional verificada y trazable que represente $\kx$, estime el campo de temperatura y documente la exactitud, la consistencia física y la robustez del análisis.

\subsubsection{PROBLEMAS ESPECÍFICOS}

Los problemas específicos descomponen la brecha general en cuatro necesidades relacionadas: especificar la representación física del objeto, disponer de referencias de verificación, evaluar escenarios heterogéneos controlados y determinar su efecto sobre la ampacidad. Esta articulación evita tratar cada problema de forma aislada y establece la secuencia lógica sintetizada en la Figura~\ref{fig:articulacion-problemas-especificos}. La sistematización posterior de la evidencia y la formulación de principios de diseño no constituyen un problema específico adicional; corresponden a las actividades de evaluación, síntesis y comunicación propias de DSR, descritas en el Capítulo~\ref{ch:metodologia}.

\begin{enumerate}[label=PE\arabic*.]
  \item La representación física requerida por el artefacto no está formalizada en términos de entradas, requisitos funcionales y criterios de calidad que ordenen el cable multicapa, las interfaces, las fuentes, los contornos y $\kx$ en escenarios heterogéneos.
  \item No se dispone de un paquete reproducible de soluciones analíticas o manufacturadas, referencias FEM con convergencia y casos publicados que permita construir y verificar el artefacto dentro del dominio delimitado.
  \item No se ha cuantificado de forma trazable cómo el contraste, el patrón, la extensión y la proximidad de las heterogeneidades modifican $\Tmax$, el campo $T(x,y)$ y la localización del punto caliente.
  \item No se han identificado las condiciones que producen la mayor discrepancia de ampacidad entre una representación 2D heterogénea verificada y una representación homogénea equivalente.
\end{enumerate}

\begin{figure}[htbp]
\centering
\resizebox{\textwidth}{!}{%
\begin{tikzpicture}[
  >=Latex,
  problema/.style={draw, rounded corners=3pt, minimum height=2.35cm, text width=4.15cm, align=center, font=\small, text=black, inner sep=6pt},
  insumo/.style={align=center, font=\small, text width=4.15cm},
  flujo/.style={-{Stealth[length=3mm]}, line width=1.1pt},
  dato/.style={font=\small\bfseries, fill=white, inner sep=1.5pt}
]
  \node[problema, fill=lightgray] (pe1) {\textbf{PE1}\\Representación física no formalizada};
  \node[problema, fill=lightblue, right=1.25cm of pe1] (pe2) {\textbf{PE2}\\Referencias de verificación no integradas};
  \node[problema, fill=green!10, right=1.25cm of pe2] (pe3) {\textbf{PE3}\\Evaluación de escenarios no trazable};
  \node[problema, fill=yellow!15, below=2.25cm of pe2] (pe4) {\textbf{PE4}\\Discrepancias de ampacidad no identificadas};

  \node[insumo, above=0.55cm of pe1] (i1) {Cable multicapa, interfaces, fuentes, contornos y $\kx$};
  \node[insumo, above=0.55cm of pe2] (i2) {Soluciones analíticas, manufacturadas, FEM y casos publicados};
  \node[insumo, above=0.55cm of pe3] (i3) {Contraste, patrón, extensión y proximidad};
  \node[insumo, below=0.55cm of pe4] (i4) {Representaciones homogénea y heterogénea 2D};

  \draw[-{Stealth[length=3mm]}, line width=1.1pt] (i1) -- (pe1);
  \draw[-{Stealth[length=3mm]}, line width=1.1pt] (i2) -- (pe2);
  \draw[-{Stealth[length=3mm]}, line width=1.1pt] (i3) -- (pe3);
  \draw[-{Stealth[length=3mm]}, line width=1.1pt] (i4) -- (pe4);

  \draw[flujo] (pe1) -- node[dato] {D1} (pe2);
  \draw[flujo] (pe2) -- node[dato] {D2} (pe3);
  \draw[flujo] (pe3.south) -- node[dato] {D3} (pe4.north);

  \node[problema, fill=red!10, right=1.25cm of pe4] (salida) {\textbf{Resultado lógico}\\Base para formular los objetivos específicos};
  \draw[flujo] (pe4) -- (salida);

\end{tikzpicture}%
}
\caption[Articulación de los problemas específicos]{Articulación secuencial de los cuatro problemas específicos de la investigación. Cada problema recibe los elementos indicados y produce una definición o evidencia necesaria para abordar el siguiente. La síntesis metodológica posterior pertenece al ciclo DSR y no se formula como un problema específico adicional. Fuente: Elaboración propia.}
\label{fig:articulacion-problemas-especificos}
\end{figure}

\subsection{JUSTIFICACIÓN Y ALCANCES}

\subsubsection{JUSTIFICACIÓN}

\paragraph{Justificación práctica.}

La justificación práctica responde al problema real: reducir la incertidumbre que puede conducir a decisiones inseguras o excesivamente conservadoras sobre la ampacidad. La investigación permitirá identificar configuraciones donde un valor homogéneo puede ocultar una restricción térmica local. Esta posibilidad se relaciona con los cambios de capacidad por humedad y resistividad reportados por \textcite[1,6--9]{khumalo2025}.

Además, el estudio puede apoyar decisiones de diseño, verificación y operación. Esto aplica cuando la evacuación de energía depende de circuitos enterrados con bajo margen térmico \parencite[14842,14851]{fariz2026}.

\paragraph{Justificación tecnológica.}

La justificación tecnológica corresponde a la respuesta propuesta para el problema técnico. El artefacto combinará conocimiento físico y aprendizaje automático para producir un campo térmico diferenciable, ya que incorporará la PDE en la función de pérdida durante el entrenamiento \parencite[686--690]{raissi2019}. Además, buscará reducir la dependencia exclusiva de pares etiquetados generados por FEM, la cual aparece en enfoques híbridos recientes para cables \parencite[1--3,12]{aldulaimi2024}.

Sin embargo, la propuesta no pretende reemplazar a FEM ni a las normas IEC. En cambio, busca establecer dentro de qué condiciones una PINN alcanza exactitud, consistencia física y robustez suficientes \parencite[93--99]{cigre2025}.

\paragraph{Justificación metodológica.}

El enfoque DSR hace investigable el problema técnico al vincular la situación física con la construcción y evaluación del artefacto \parencite[55--56]{peffers2007}. Por tanto, el estudio no se limita a ejecutar software, sino que aportará un procedimiento reproducible de especificación y verificación. También propondrá principios de diseño sobre interfaces y heterogeneidad del terreno. Este énfasis sigue la distinción entre artefacto situado y contribución abstracta de \textcite[342--350]{gregor2013}.

\paragraph{Justificación científica y académica.}

La justificación científica responde al problema subyacente: producir conocimiento sobre cuándo y en qué magnitud la heterogeneidad espacial modifica la temperatura y la ampacidad frente a una representación homogénea equivalente. El trabajo conectará dos cuerpos de conocimiento: el dimensionamiento térmico de cables enterrados y el aprendizaje informado por física.

Los resultados de PINN en dominios térmicos controlados ya han sido reportados \parencites[1--3,17]{xing2023}[1,3,16]{pan2025}. Sin embargo, aquí se evaluarán en una geometría multimaterial, donde el entrenamiento y las interfaces todavía deben verificarse \parencite[2--6]{ren2025}.

\subsubsection{ALCANCES Y LIMITACIONES}

\paragraph{Alcances.}

El alcance temático comprende la conducción de calor en secciones bidimensionales de cables XLPE enterrados, considerando la heterogeneidad del suelo nativo, el lecho de asiento y el suelo de relleno. El análisis calculará la temperatura y la ampacidad en estas condiciones.

El alcance metodológico comprende el diseño y la evaluación de un artefacto basado en una PINN mediante la ciencia del diseño (DSR). Sus resultados se compararán con referencias analíticas, soluciones obtenidas mediante el método de elementos finitos (FEM) y casos publicados.

El alcance espacial considera configuraciones representativas de circuitos con cables enterrados para evacuar la energía de centrales solares y eólicas, sin limitarse a una instalación propietaria. El estudio se desarrollará entre julio de 2026 y junio de 2027, y la revisión bibliográfica tendrá como fecha de corte junio de 2026.

El producto incluirá la arquitectura lógica, el prototipo de la PINN, los datos de los escenarios, las pruebas y el informe de evaluación. También comprenderá principios de diseño reproducibles que faciliten la verificación del artefacto.

\paragraph{Limitaciones.}

El estudio será principalmente estacionario y bidimensional, no incluye el régimen transitorio. El modelo no resolverá la interacción electromagnética completa del cable. Por tanto, las pérdidas eléctricas se representarán como una fuente de calor por efecto Joule, proporcional a $I^2R$ \parencites[752--760]{neher1957}[cláusulas 4--5]{iec60287}.

Tampoco se modelarán envejecimiento dieléctrico, humedad acoplada, migración de humedad ni convección subterránea. En consecuencia, las propiedades del suelo nativo, del lecho de asiento y del suelo de relleno se tratarán como atributos físicos de cada realización de $E_3$, no como VI metodológicas adicionales.

La evaluación del artefacto se limitará a verificación y comparación documentada. Por verificación se entenderá la comparación con casos analíticos, soluciones manufacturadas, referencias FEM y casos publicados \parencites{cigre2022}[3--4,93--99]{cigre2025}.

No se realizará validación experimental de campo, debido a que requeriría instrumentación y seguimiento de una instalación real. Por tanto, los resultados no se generalizarán fuera de las geometrías, materiales, fuentes y contornos evaluados.

Además, cualquier ventaja computacional se reportará junto con hardware, tolerancias, semillas e inicializaciones. Así se comunicará la evidencia con límites claros, como exige la metodología DSR \parencite[350--351]{gregor2013}.

El estudio no incluirá una evaluación económica de los cables o de la instalación; por tanto, no estimará CAPEX, OPEX, costos de construcción ni retorno de inversión. Toda referencia a desempeño o recursos corresponderá exclusivamente al ámbito computacional y se reportará como información secundaria de ejecución y reproducibilidad.

